\section{What is a quiver?}\label{sec:12-path-algebras:01-what-is-a-quiver:1}

The stated goal of this chapter is the path algebra of a directed graph, a construction that needs paths before it needs an algebra, and needs a graph before it needs paths. So the chapter starts at the bottom, the graph itself, stripped to the bare minimum of data a ``directed graph'' requires. Everything from here through the checkpoint project is built on top of exactly this one definition.

A path is a sequence of arrows chained end to end, each one starting where the last one ended. Stating that requirement precisely, ``where an arrow starts'' and ``where it ends,'' is already the whole definition needed: a set of things arrows can start and end at, call them vertices, and a set of arrows, each one tagged with which vertex it starts at and which it ends at. Nothing else is required to make sense of chaining two arrows together, so nothing else goes into the definition.

A \textbf{quiver} is exactly this: a set of vertices \(Q_0\), a set of arrows \(Q_1\), and two functions \(s, t : Q_1 \to Q_0\) recording the \textbf{source} and \textbf{target} of each arrow. It is a \emph{directed graph} under another name. This is the same notion as \href{https://loogle.lean-lang.org/?q=Quiver}{\texttt{Quiver}} in Mathlib (built here from scratch, following the ``no Mathlib'' policy of Chapter 1, rather than reusing the Mathlib version).

Picture an arrow \(\alpha : i \to j\) as a literal arrow drawn from vertex \(i\) to vertex \(j\); \(s(\alpha) = i\) and \(t(\alpha) = j\).

\subsection{Example quiver}\label{sec:12-path-algebras:01-what-is-a-quiver:2}

Take vertices \(\{1, 2, 3\}\) and arrows

\[
\alpha : 1 \to 2, \qquad \beta : 2 \to 3
\]

\subsection{Sources, quoted}\label{sec:12-path-algebras:01-what-is-a-quiver:3}

Formal definitions and citations for this section, gathered here for reference (full entries in the \hyperref[sec:root:bibliography:1]{Bibliography}):

\begin{itemize}
\tightlist
\item
  \textbf{Quiver.} Stated verbatim as \(Q = (Q_0, Q_1, s, t)\) (\cite{AssemSimsonSkowronski2006}, Definition 1.1, Ch. II §1 ``Quivers and path algebras''), a set of vertices \(Q_0\), a set of arrows \(Q_1\), and two functions \(s, t : Q_1 \to Q_0\) giving the source and target of each arrow.
\item
  Schiffler (\cite{Schiffler2014}) covers an elementary, textbook-level treatment of the same definition.
\end{itemize}
